Clinical narratives in electronic health records (EHRs) are indispensable for understanding patient status, informing decision support, and enabling biomedical research. Narrative notes—spanning discharge summaries, progress notes, imaging and operative reports—capture observations, differential diagnoses, treatments, test results, and clinician reasoning that structured fields rarely encode. Expressions such as ``intermittent chest pain worsening with exertion but relieved by rest'' or ``likely drug-induced hepatitis, possibly from isoniazid'' capture nuance that structured fields miss; measurements like cardiac ejection fraction values or tumor size often do not fit neatly into predefined problem lists and billing codes yet are central to clinical interpretation~\new{\cite{Tayefi_2021,Sushil_2024}}. Interpretation is not just an important task of clinical interpretation for immediate clinician decisions. It also is the translation of this text to codified terminologies (e.g., CPT codes) or numerical representations (e.g., lab result tables) to drive tasks in epidemiology, automated decision support, reimbursement, and trial eligibility, among other downstream uses of electronic healthcare data.

Not only are the structured EHR entries incomplete or inaccurate, but conventional natural language processing methods—such as standard named entity recognition—often fail to capture the nuance of clinical reasoning, contextual qualifiers, or temporal relationships in notes \cite{Kohane_2021}. Even when structured codes or extracted entities are available, they may not reflect whether a condition was truly present, a medication was taken, or a disease state had recurred \cite{Kim_2023}. Consequently, high-quality labels in practice still rely on chart review: investigators read notes to adjudicate diagnoses and verify supporting evidence (medications, symptoms, tests). Yet, chart review is slow, costly, and variable across annotators; double abstraction is frequently required to ensure quality. This persistent reliance highlights the need for automated, scalable approaches that can more effectively extract, structure, and standardize clinical information consistently while leaving an auditable trace.

Classical approaches to clinical natural language processing have attempted to automate extraction and structuring of information from narrative notes but face important limitations. Rule- and lexicon-driven systems, such as cTAKES \cite{Savova_2010}, MetaMap \cite{Aronson2001MetaMap}, and MedLEE \cite{Friedman_1995}, map entities to controlled vocabularies (e.g., the Unified Medical Language System), using handcrafted rules and dictionaries for concept recognition, negation detection, and temporal normalization. While these pipelines established the foundation for clinical NLP, they struggle with shorthand, spelling variants and errors, as well as non-standard grammar common in clinical notes, and they require substantial manual engineering to adapt across institutions and use cases \cite{Tayefi_2021,Pathak_2013}. \new{Between these rule-based systems and modern large language models, a generation of supervised neural methods advanced clinical entity recognition and normalization---BiLSTM-CRF sequence taggers and, subsequently, domain-pretrained transformer encoders such as BioBERT and ClinicalBERT~\cite{Lample_2016,Lee_2020,Alsentzer_2019}. These improved span detection but require sizeable task-specific labeled corpora, remain brittle to phrasing variation and institutional shift, and typically address entity detection in isolation rather than joint normalization, attribute, and temporal extraction.}

Building on these limitations of rule- and lexicon-driven systems, transformer-based large language models—such as GPT-4 \cite{https://doi.org/10.48550/arxiv.2303.08774} and Med-PaLM \cite{Singhal_2023}—have reshaped NLP and show promise for clinical tasks including entity and relation extraction, summarization, and question answering \cite{https://doi.org/10.48550/arxiv.2005.14165}. Domain-scaled models, such as GatorTron \cite{https://doi.org/10.48550/arxiv.2203.03540} and Clinical-MobileBERT \cite{Rohanian_2023,Rohanian_2024}, report strong performance on benchmarks such as i2b2 concept extraction \new{\cite{Uzuner_2011}}, and instruction-tuned assistants can generalize in zero- and few-shot regimes \cite{Radford2018Improving}. Yet, important barriers remain for safe clinical deployment extraction: ontology grounding to UMLS/SNOMED CT is often imprecise with occasional hallucinated codes \cite{Karabacak_2023,https://doi.org/10.48550/arxiv.2309.05922}; context and output-length limits complicate reasoning across long clinical notes \cite{https://doi.org/10.48550/arxiv.2303.18223,https://doi.org/10.48550/arxiv.2407.21783}; and in real-world datasets (e.g., from the 4CE consortium), note lengths frequently exceed such limits, hindering cross-sectional reasoning \cite{Brat_2020}. As a result, single-pass prompting remains unreliable for chart-review-like decisions at scale.

To address these gaps, we present CLINES—a \new{Clinical LLM-based Information Extraction and Structuring} agent that brings the consistency of chart review into an automated pipeline. Unlike prior systems limited by handcrafted rules or naive single-pass prompting, CLINES integrates several innovations: informed segmentation for processing long notes; reasoning-capable LLM extraction of entities and relations; attribution of clinical context (e.g., negation, experiencer, dosage, laboratory values); normalization to existing ontologies (e.g., UMLS); temporal alignment of events; and reconciliation into an i2b2-style schema to facilitate downstream integration. By uniting these components, CLINES represents a substantial step forward in clinical information extraction, with improvements over both classical NLP pipelines and baseline LLM prompting. Its improvements are validated using real-world EHR notes spanning multiple disease areas and institutions, \new{underscoring accuracy across the disease areas and institutions evaluated; we nonetheless temper claims of broad generalizability, as several common note types (e.g., discharge summaries, operative and radiology reports) remain unevaluated}. Equally important, CLINES is engineered to enable different generative models to be substituted within the workflow as dictated by user or institutional preferences regarding cost, property, performance and security. \new{CLINES assumes de-identified, plain-text notes as input; upstream de-identification, format normalization, and document segmentation lie outside the present scope and are addressed as deployment considerations in the Discussion.}
